\IEEEPARstart{S}{imBench} targets simulator-grade digital-twin (DT) generation for Project Chrono/PyChrono, and spans a diverse set of scenarios drawn from multibody systems (MBS), finite element analysis (FEA), vehicle dynamics (VEH), sensor integration (SEN), and robotics (RBT). 
Across these categories, SimBench evaluates an S-LLM’s ability to (i) construct simulator-ready DT artifacts from natural-language requests, and (ii) iteratively refine and extend existing DT implementations in a multi-turn workflow (including targeted modifications and bug fixing). 
In this work, S-LLMs are treated as black-box code generators; developing or selecting S-LLMs is outside the scope. Our focus is the benchmark design and the judge model (J-LLM) used for interpretable evaluation.

\vspace{2pt}
\noindent The families of DTs assessed in SimBench include:
\begin{itemize}
    \item \textbf{Multibody Systems (MBS):} Mechanism and contact-rich systems (e.g., articulated arms, gear trains, slider--crank mechanisms) that exercise rigid-body dynamics, constraints, and contact/friction modeling.
    \item \textbf{Finite Element Analysis (FEA):} Deformable-body scenarios (e.g., cables, beams, shells, plates) that test proficiency in discretization choices, boundary conditions, loading, and solver configuration.
    \item \textbf{Vehicle Dynamics (VEH):} On-road and off-road vehicles (e.g., city bus, HMMWV, M113, Kraz, MAN, sedans) with driver/powertrain/tire models and scenario logic (e.g., lane changes, logging, mixed traffic).
    \item \textbf{Sensor Integration (SEN):} GPS/IMU/LiDAR/camera configurations that evaluate correct sensor instantiation, attachment, calibration, and data acquisition pipelines.
    \item \textbf{Robotics (RBT):} Robotic platforms (e.g., Turtlebot, Curiosity, VIPER) and off-road interactions such as deformable terrain and soil contact models (SCM) \cite{Krenn2011}.
\end{itemize}
SimBench draws on $N_{\text{sys}}=34$ expert-authored physical systems, each realized as a three-turn demonstration, yielding $N_{\text{task}}=102$ turn-level benchmark tasks across these categories. Each system is represented by three high-quality turns with progressively increasing complexity, enabling robust assessment of code generation and code modification capabilities. Figure~\ref{fig:demo_gallery} illustrates representative SimBench scenarios.

\begin{figure}[h]
    \centering
    \includegraphics[width=0.5\textwidth]{demo_overview.pdf}
    \caption{A subset of simulation scenarios in SimBench.}
    \vspace{-10pt}
    \label{fig:demo_gallery}
\end{figure}

\medskip

{\noindentBold{Task Formulation and Multi-Turn Protocol}}: 
SimBench tasks are expressed as natural-language requests that vary in specificity. We distinguish two request styles:

\begin{itemize}
    \item \textbf{Vague Requests:} high-level instructions that intentionally leave parts of the DT specification underspecified (e.g., missing parameter values or implementation details). These prompts test whether an S-LLM can infer reasonable defaults, make meaningful modeling choices, and produce a coherent simulator-ready DT under incomplete information.
    \item \textbf{Sharp Requests:} explicit, constraint-rich instructions that precisely specify the desired DT changes and expected behaviors. These prompts test faithful instruction following and targeted code editing in response to concrete requirements.
\end{itemize}

Each SimBench task is structured as a \emph{three-turn} interaction that reflects a common DT-development workflow: an initial high-level request followed by increasingly specific refinements. In our benchmark instantiation, the first turn is \emph{typically} a Vague Request, while the second and third turns are Sharp Requests.
Table~\ref{tab:mt_bench_sample} provides representative multi-turn prompts across categories.

\begin{enumerate}
    \item \textbf{Turn 1 (Vague Request): DT construction from scratch.}
    The first turn requests building a DT from an abstract or minimally specified description (e.g., a basic mechanism, a deformable object setup, or a vehicle/robot scenario). The goal is to probe foundational knowledge of PyChrono primitives and the ability to make principled modeling decisions when details are missing.

    \item \textbf{Turn 2 (Sharp Request): targeted modification and bug fixing.}
    The second turn provides concrete requirements for modifying the previously generated DT (e.g., replacing a joint type, adding a sensor, changing initial conditions, or adjusting solver/visualization settings). In addition, SimBench monitors \emph{bug-fixing} requests where a DT implementation is intentionally perturbed (e.g., API misuse, inconsistent parameters, or missing configuration), and the S-LLM is instructed to identify and correct the code.

    \item \textbf{Turn 3 (Sharp Request): challenging extension of the scenario.}
    The final turn specifies a more demanding extension that typically requires multi-step reasoning and coordinated changes across components (e.g., introducing autonomy/control logic, executing a lane-change scenario, integrating multiple sensors, adding additional agents/objects, or logging task-relevant state). This turn evaluates an S-LLM’s ability to maintain consistency with prior turns while scaling the DT to a more complex virtual experiment.
\end{enumerate}

\begin{table}[h!]
	\centering
	\begin{tabular}{|c|c|c|c|}
		\hline
		\multirow{2}{0.7cm}{Genre} & \multicolumn{3}{c|}{Multi-turn Task Description}  \\ \cline{2-4}
		& 1st turn & 2nd turn & 3rd turn   \\ \hline
		\parbox[c][1.75cm][c]{0.7cm}{MBS} & \parbox[c][1.75cm][c]{1.5cm}{Generate a one-link arm} & \parbox[c][1.75cm][c]{2.1cm}{Replace revolute joint with spherical joint} & \parbox[c][1.75cm][c]{2.1cm}{Extend single link into double-link arm}  \\ \hline
		\parbox[c][2cm][c]{0.7cm}{FEA} & \parbox[c][2cm][c]{1.5cm}{Model the folding of tablecloth} & \parbox[c][2cm][c]{2.1cm}{Change viz. \& solver \& apply load; set boundary conditions} & \parbox[c][2cm][c]{2.1cm}{Enhance to higher quality visualization}  \\ \hline
		\parbox[c][2.5cm][c]{0.7cm}{VEH} & \parbox[c][2.5cm][c]{1.5cm}{Create Kraz truck vehicle model} & \parbox[c][2.5cm][c]{2.1cm}{Adjust truck initial position and viz. angle; implement double lane change} & \parbox[c][2.5cm][c]{2.1cm}{Switch scene to highway; add another sedan; record truck state info}  \\ \hline
		\parbox[c][2cm][c]{0.7cm}{SEN} & \parbox[c][2cm][c]{1.5cm}{Generate LiDAR to sense obj. defined via mesh} & \parbox[c][2cm][c]{2.1cm}{Replace mesh with box object and add one more 2D LiDAR sensor} & \parbox[c][2cm][c]{2.1cm}{Create LiDAR and camera sensors attaching to a scaled vehicle}  \\ \hline
		\parbox[c][1.75cm][c]{0.7cm}{RBT} & \parbox[c][1.75cm][c]{1.5cm}{Generate a Turtlebot robot} & \parbox[c][1.75cm][c]{2.1cm}{Showcase robot moving right, left, and straight ahead} & \parbox[c][1.75cm][c]{2.1cm}{Add 5 boxes randomly \& use LiDAR to monitor scene}  \\ \hline
	\end{tabular}
	\caption{Sample task requests submitted to the S-LLM.}
	\vspace{-10pt}
	\label{tab:mt_bench_sample}
\end{table}


\medskip
\noindent\textbf{Example: Curiosity rover (RBT) multi-turn DT generation (abridged).}
\label{app:curiosity_example}
We provide an illustrative robotics (RBT) task to demonstrate the three-turn protocol.
Turn~1 is a \emph{vague} request that specifies the overall DT goal, while Turns~2--3 are \emph{sharp} requests that require targeted code edits and extensions.
In SimBench, Turns~2--3 additionally include the full code context from the previous turn; we omit the code context here for readability.

\begin{PromptBlock}{Curiosity (RBT) --- Turn 1 (Vague Request; DT construction)}
    Create a PyChrono simulation showcasing a Curiosity rover navigating on a rigid terrain.
    The system includes a fixed ground with collision properties and a rover model equipped with a motor control driver for real-time steering inputs.
    Visualize the simulation using the Irrlicht visualization system with customizable camera settings, lighting, shadows, and additional visual elements like textures and logos.
    \end{PromptBlock}
    
    \begin{PromptBlock}{Curiosity (RBT) --- Turn 2 (Sharp Request; code editing)}
    Do the following changes based on the provided code:
    (*@\textbf{1.}@*) Add a long box obstacle for the rover to cross.\\
    (*@\textbf{2.}@*) Change the initial position of the rover from $(0,\,0.2,\,0)$ to $(-5,\,0.0,\,0)$.\\
    (*@\textbf{3.}@*) Set the rover to move forward with zero steering input, instead of gradually increasing steering input over time.
    \end{PromptBlock}
    
    \begin{PromptBlock}{Curiosity (RBT) --- Turn 3 (Sharp Request; sensor integration)}
    Do the following changes to the given code:
    (*@\textbf{1.}@*) Add the sensor module (\texttt{pychrono.sensor}) import.\\
    (*@\textbf{2.}@*) Create a sensor manager (\texttt{sens.ChSensorManager(system)}).\\
    (*@\textbf{3.}@*) Add a LiDAR sensor attached to the rover chassis:
    \begin{itemize}
      \item Define LiDAR parameters such as update rate, horizontal and vertical samples, and field of view.
      \item Create the LiDAR sensor (\texttt{sens.ChLidarSensor}) and attach it to the rover chassis.
      \item Set up filters for the LiDAR data.
      \item Add the LiDAR sensor to the sensor manager.
    \end{itemize}
    (*@\textbf{4.}@*) Update the sensor manager in the simulation loop (\texttt{manager.Update()}).
    \end{PromptBlock}
    
    
% Optional (recommended): one sentence to reinforce "interpretability" link
\noindent This example also motivates the need for interpretable evaluation: the sharp requests in Turns~2--3 require localized, component-level modifications (obstacle geometry, initial conditions, driver logic, and sensor pipeline), which are difficult to diagnose using a single pass/fail outcome.

\medskip

{\noindentBold{Evaluation framework and modalities}:}
\label{sec:modalities} In the approach proposed, a rubric-based J-LLM evaluates the performance of an S-LLM using the tasks contained in the SimBench environment. Unlike traditional metrics such as BLEU, ROUGE, and CodeBLEU, the J-LLM is designed to provide a comprehensive evaluation of the S-LLM's performance in multi-turn simulation tasks. The J-LLM is instructed via a rubric-driven prompt that provides (i) the candidate code, 
(ii) an expert reference solution (when available), and optionally (iii) a compact API documentation excerpt.
The prompt elicits both a scalar score and itemized, component-level deductions.

Three evaluation modalities are considered in SimBench.
In \textbf{J-LLM\_Ref\_Doc}, the judge is prompted with both the expert reference code and an expert-written API documentation excerpt (approximately 4{,}000 tokens), in addition to the S-LLM’s candidate code.
This setting requires a long context window to jointly accommodate documentation, reference, and candidate code; in our experiments, this constraint limits the applicability of judge models with shorter context windows.

In \textbf{J-LLM\_Ref}, the judge is provided the expert reference code and the candidate code, without additional documentation, and must rely on general coding priors together with the reference implementation.

Finally, in \textbf{J-LLM\_Doc}, the judge is provided only the expert-written API documentation excerpt and the candidate code, without access to the reference implementation

\medskip

{\noindentBold{Execution-Based Evaluation}}: In addition to J-LLM evaluation, SimBench incorporates execution-based metrics that provide ground-truth validation of generated code:

\begin{itemize}
    \item \textbf{Compile@1}: Measures whether the generated code is parsed successfully in Python. This metric captures the S-LLM's ability to produce syntactically correct PyChrono code. For simulators other than PyChrono, this test might amount to a C or C++ compilation.
    \item \textbf{Pass@1}: Measures whether the generated code passes predefined functional tests. Unlike Compile@1, Pass@1 requires the code to be semantically correct and to produce the expected simulation behavior, as evidenced by a successful pass, correct numerical outputs, and physically reasonable visualizations.
\end{itemize}

\medskip

{\noindentBold{Code Similarity Metrics}}: SimBench also computes automated code similarity metrics:

\begin{itemize}
    \item \textbf{CodeBLEU}~\cite{ren2020codebleu}: A composite metric that considers n-gram match, weighted n-gram match, syntax match (AST), and dataflow match. CodeBLEU captures both surface-level and structural code similarity.
    \item \textbf{ROUGE-L}: Measures the longest common subsequence between generated and reference code, capturing sequential similarity.
\end{itemize}

\medskip
\noindent\textbf{Rubric-prompt design for interpretable judging.}
SimBench adopts a rubric-driven judge prompting strategy to obtain both a scalar score and localized, component-level feedback.
The J-LLM is prompted with (i) the candidate DT code produced by an S-LLM, (ii) an expert reference implementation for the same turn, and optionally (iii) PyChrono API documentation.
The prompt enforces a fixed, additive deduction scheme across six categories (Completeness, Correctness, Code Quality, Efficiency, Robustness, and Visualization), with explicit point-deduction guidelines.
This structure encourages the judge to (a) identify \emph{what} is missing or incorrect, (b) explain \emph{why} it matters relative to the reference or best practices, and (c) translate that diagnosis into a consistent numeric score.

\vspace{2pt}
\noindent
The rubric categories are intentionally aligned with common failure modes in simulator-grade DT code generation.
For example, structural omissions (e.g., missing sensor manager updates, missing collision settings, absent visualization blocks) map naturally to \emph{Completeness}; API misuse and frame/pose inconsistencies map to \emph{Correctness}; and brittle or non-idiomatic implementations map to \emph{Robustness} and \emph{Code Quality}.
To mitigate known LLM-judge biases, the prompt additionally includes guardrails against position bias, length bias, and stylistic preferences. It also requires the final score to be emitted in a machine-parsable format \texttt{[[x]]} for downstream aggregation.

\vspace{2pt}
\noindent\textbf{Judge calibration protocol.}
To calibrate the rubric-grounded judge prompt, we constructed a calibration set covering the major Chrono task classes. Concretely, we selected one representative task from each of five classes (\texttt{MBS}, \texttt{FEA}, \texttt{SEN}, \texttt{RBT}, \texttt{VEH}), yielding $N_{\text{ref}}=5$ expert reference scripts. We use one task per class because tasks within a class share core APIs, structure, and recurring failure modes, making a representative reference sufficient for prompt calibration without overfitting to a narrow instance.

For each reference script, we manually synthesized $N_{\text{var}}=5$ perturbed variants by injecting common implementation errors frequently discussed in the Project Chrono community forum \cite{projectChronoForum}, while keeping the underlying task intent unchanged. The injected bugs span: (i) API hallucination or misuse (incorrect class/method names, deprecated calls, invalid signatures), (ii) constraint/joint specification errors (missing constraints, incorrect frames, inconsistent coordinates), (iii) simulation-loop and subsystem omissions (e.g., missing sensor-manager updates, missing module initialization, incorrect update order), (iv) integrator/solver misconfiguration (inappropriate system type, solver choice, time step, or tolerances), and (v) visualization/IO omissions (missing visualization setup, assets, render/update calls, or output logging).

We iteratively refined the judge prompt using only this calibration set until (a) the induced ranking of a small panel of S-LLMs agreed with expert rankings, and (b) no systematic judge failure modes were observed on a small holdout subset of perturbed cases.

\medskip

\noindent\textbf{J-LLM scoring, aggregation, and reporting.}
SimBench contains $N_{\text{sys}}=34$ distinct physical systems. Each system is evaluated through a fixed three-turn protocol ($t\in\{1,2,3\}$), yielding $N_{\text{task}}=3N_{\text{sys}}=102$ turn-level benchmark tasks.

For a given S-LLM, let $y_{i,t}$ denote the code produced for system $i$ at turn $t$. Under a chosen evaluation modality (Sec.~\ref{sec:modalities}), the judge model (J-LLM) is \emph{prompted} to output (i) an interpretable textual critique with itemized deductions and (ii) a scalar score:
\begin{equation}
s_{i,t} \;\triangleq\; \mathrm{JLLM}\!\left(y_{i,t};\, r_{i,t}, d\right)\in[0,100],
\end{equation}
where $r_{i,t}$ is the expert reference implementation when available, and $d$ denotes the optional API-documentation context.

\vspace{2pt}
\noindent\emph{System-level aggregation.}
We first aggregate the three turn scores into a system-level score:
\begin{equation}
S_i \;\triangleq\; \sum_{t=1}^{3} w_t\, s_{i,t},
\qquad \text{with}\quad \sum_{t=1}^{3} w_t = 1,
\label{eq:sys_agg}
\end{equation}
using uniform weights $w_1=w_2=w_3=\tfrac{1}{3}$ to reflect end-to-end workflow competence.
\vspace{2pt}

\noindent\emph{Overall score.}
We report the overall SimBench score as the task-level macro-average over all $N_{\text{task}}=102$ turn-level tasks:
\begin{equation}
S \;\triangleq\; \frac{1}{N_{\text{task}}}\sum_{i=1}^{N_{\text{sys}}}\sum_{t=1}^{3} s_{i,t}.
\label{eq:overall_agg}
\end{equation}
Equivalently, since each system contributes exactly three turns, $S = \frac{1}{N_{\text{sys}}}\sum_{i=1}^{N_{\text{sys}}} S_i$.

\vspace{2pt}
\noindent\emph{Category-level reporting.}
Let $\mathcal{T}_c$ denote the set of \emph{turn-level} tasks $(i,t)$ that belong to category $c$. We report category scores as
\begin{equation}
S^{(c)} \;\triangleq\; \frac{1}{|\mathcal{T}_c|}\sum_{(i,t)\in\mathcal{T}_c} s_{i,t}.
\label{eq:cat_agg}
\end{equation}

{\noindent\textbf{Dataset Complexity Analysis}}: To contextualize SimBench's complexity, we provide a quantitative comparison with established code generation benchmarks in Table~\ref{tab:benchmark_comparison}. Token counts are computed using OpenAI's \texttt{tiktoken} library with the \texttt{cl100k\_base} encoding.

\begin{table}[h!]
    \centering
    \caption{Comparison of SimBench with existing code generation benchmarks. SimBench exhibits significantly higher complexity in both prompt and solution lengths.}
    \label{tab:benchmark_comparison}
    \begin{tabular}{l r r r r r}
        \toprule
        \textbf{Benchmark} & \textbf{Tasks} & \textbf{Prompt} & \textbf{Solution} & \multicolumn{2}{c}{\textbf{Ratio to SimBench}} \\
        & & \textbf{(tokens)} & \textbf{(tokens)} & Prompt & Solution \\
        \midrule
        MBPP~\cite{austin2021program} & 500 & 16 & 58 & 61.6$\times$ & 24.4$\times$ \\
        MBPP+~\cite{liu2024your} & 378 & 19 & 40 & 51.8$\times$ & 35.4$\times$ \\
        DS-1000~\cite{lai2023ds} & 1,000 & 282 & 42 & 3.5$\times$ & 33.7$\times$ \\
        HumanEval~\cite{chen2021evaluating} & 164 & 131 & 54 & 7.5$\times$ & 26.2$\times$ \\
        HumanEval+~\cite{liu2024your} & 164 & 131 & 54 & 7.5$\times$ & 26.2$\times$ \\
        BigCodeBench~\cite{zhuo2024bigcodebench} & 1,140 & 145 & 112 & 6.8$\times$ & 12.6$\times$ \\
        CodeContests~\cite{li2022competition} & 165 & 593 & 762 & 1.7$\times$ & 1.9$\times$ \\
        \midrule
        \textbf{SimBench (ours)} & \textbf{102} & \textbf{985} & \textbf{1,415} & 1.0$\times$ & 1.0$\times$ \\
        \bottomrule
    \end{tabular}
\end{table}

A key distinction of SimBench lies in its multi-turn structure. Unlike single-turn benchmarks where each task is independent, SimBench's second and third turns require the S-LLM to modify existing code from the previous turn. This design reflects real-world software development workflows where developers iteratively refine and extend existing codebases. Table~\ref{tab:turn_complexity} details the prompt composition across turns.

\begin{table}[h!]
    \centering
    \caption{Prompt complexity across SimBench turns. Turns 2 and 3 include full code context from previous turns, reflecting realistic code modification scenarios.}
    \label{tab:turn_complexity}
    \begin{tabular}{l r r r r}
        \toprule
        \textbf{Turn} & \textbf{Text} & \textbf{Code Context} & \textbf{Total Prompt} & \textbf{Solution} \\
        & \textbf{(tokens)} & \textbf{(tokens)} & \textbf{(tokens)} & \textbf{(tokens)} \\
        \midrule
        Turn 1 & 87 & --- & 87 & 1,252 \\
        Turn 2 & 189 & 1,252 & 1,441 & 1,416 \\
        Turn 3 & 123 & 1,306 & 1,429 & 1,577 \\
        \midrule
        \textbf{Average} & \textbf{133} & \textbf{853} & \textbf{985} & \textbf{1,415} \\
        \bottomrule
    \end{tabular}
\end{table}

The multi-turn paradigm introduces several challenges absent in conventional benchmarks:
\begin{itemize}
    \item \textbf{Long-context reasoning}: With prompts averaging 985 tokens (and exceeding 1,400 tokens in later turns), S-LLMs must maintain coherence across extended contexts while understanding the existing codebase.
    \item \textbf{Code comprehension}: Unlike generation-only tasks, turns 2 and 3 require the S-LLM to first understand the provided code before making targeted modifications, testing both reading and writing capabilities.
    \item \textbf{Incremental complexity}: Solutions grow progressively longer (1,252 $\rightarrow$ 1,416 $\rightarrow$ 1,577 tokens), requiring the S-LLM to manage increasing code complexity while maintaining consistency with prior implementations.
\end{itemize}

Table~\ref{tab:category_complexity} provides a breakdown of complexity by simulation category, revealing that Sensor (SEN) and Finite Element Analysis (FEA) scenarios demand the most extensive code generation.

\begin{table}[h!]
    \centering
    \caption{Complexity by simulation category.}
    \label{tab:category_complexity}
    \begin{tabular}{l r r r}
        \toprule
        \textbf{Category} & \textbf{Systems} & \textbf{Avg. Prompt} & \textbf{Avg. Solution} \\
        & & \textbf{(tokens)} & \textbf{(tokens)} \\
        \midrule
        SEN & 4 & 1,296 & 1,932 \\
        FEA & 5 & 1,172 & 1,598 \\
        MBS & 5 & 1,011 & 1,415 \\
        VEH & 14 & 885 & 1,312 \\
        RBT & 6 & 640 & 922 \\
        \bottomrule
    \end{tabular}
\end{table}

\noindent\textbf{Reproducibility of prompts and qualitative examples.}
All wrapper prompts used to query S-LLMs, as well as the rubric prompts used by the J-LLM, are provided in Appendix~\ref{app:prompt_templates} and Appendix~\ref{app:jllm_prompts}.
A complete three-turn Curiosity rover case study (including expert reference solutions, turn-to-turn patches, and an example S-LLM output with J-LLM feedback) is included in Appendix~\ref{app:curiosity_gt} and Appendix~\ref{app:claude_curiosity_turn3}.

\bigskip

\noindent\textbf{Generality and Portability to Other Simulators.}
Although SimBench is instantiated in this work for Chrono/PyChrono, the benchmark methodology is intended to be \emph{simulator-agnostic}. In particular, SimBench separates (i) the \emph{task protocol} used to elicit multi-turn digital-twin (DT) development behavior from an S-LLM, (ii) the \emph{context channels} used to condition evaluation (expert references and documentation), and (iii) the \emph{execution oracle} used to anchor correctness. Porting SimBench to another simulator therefore amounts to providing simulator-specific implementations of the following three interface points.

\textbf{(1) Task specification (multi-turn protocol and artifact expectations).}
SimBench tasks are defined as a fixed-turn interaction that reflects realistic DT authoring: an initial DT creation request (Turn 1), followed by successive DT modification/extension requests (Turns 2--3) that require understanding and editing the existing artifact. This protocol is independent of the simulator itself; however, the \emph{artifact type} and \emph{acceptance criteria} must be specified per target platform (e.g., Python scripts, input decks, XML/URDF/SDF configurations, plugin code, or mixed pipelines). The same prompt-genres used in SimBench (e.g., \emph{vague} vs.\ \emph{sharp} requests) can be retained to probe different aspects of instruction following, default inference, and code-editing competence.

\textbf{(2) Context sources (documentation and expert references).}
SimBench evaluates generated DTs under optional context channels: (i) simulator/API documentation and (ii) expert reference implementations when available. To port the benchmark, one should provide analogous context for the target simulator: stable API/documentation excerpts that cover the exercised features, and a small expert-authored dataset of reference DT solutions and turn-to-turn patches for the selected scenario families. These references serve two roles: they provide grounding for evaluation and they enable calibration of the rubric-based judge model (J-LLM) to match expert scoring behavior in the target ecosystem.

\textbf{(3) Execution oracle (what ``compile'' and ``pass'' mean).}
Finally, SimBench anchors evaluation using execution-based signals (e.g., Compile@1 and Pass@1), while complementing them with interpretable rubric-based judging. For a new simulator, the benchmark designer must define the platform-appropriate execution oracle, including: (i) what constitutes a syntactic validity check (``compile''), (ii) what constitutes a semantic correctness check (``pass''), and (iii) tolerances, timeouts, and determinism controls. Importantly, the oracle need not mirror PyChrono's toolchain; it must only provide objective, automated checks that capture whether the generated DT artifact is runnable and functionally correct for the intended scenario.

\textbf{Examples.}
To illustrate how these interface points specialize across platforms:
(i) For \emph{Abaqus} \cite{ABAQUSreference}, the ``compile'' analogue is typically an \texttt{abaqus job=} submission that completes without preprocessing/runtime errors (input-deck validity, successful analysis completion), while ``pass'' can be defined via regression checks on output database (ODB) quantities such as displacements/forces/energies, convergence status, and invariants under prescribed loading and boundary conditions.
(ii) For \emph{MuJoCo} \cite{todorovMujoco2012}, the syntactic check corresponds to successful model parsing/loading (XML validity, asset resolution, no warnings/errors that prevent simulation), while ``pass'' can be defined via rollout-based checks (no NaNs/instabilities, constraint satisfaction within tolerance) and task-specific quantitative comparisons (trajectory tracking error, contact event timing, sensor outputs) against reference rollouts.
(iii) For other ecosystems (e.g., OpenFOAM \cite{openFoam}, PyBullet \cite{Pybullet2023}, Isaac Sim \cite{IssacNvidia2023}, Gazebo \cite{gazebo}), the same pattern applies: the task protocol and rubric remain stable, while the artifact formats, documentation corpus, and execution oracle are adapted to the platform's compilation/loading and correctness notions.

From a high vantage point, Chrono serves as an end-to-end \emph{existence proof} for constructing a calibrated, interpretable, execution-anchored benchmark for simulator-grade DT generation. The methodology transfers to other simulators by re-instantiating the three interface points above: task specification, context sources, and execution oracle.
